% ===========================================================================
%  E05 — A experiência do estudante na plataforma
% ===========================================================================
\chapter{E05 \textperiodcentered{} A experiência do estudante na plataforma}
\label{ep:e05}

\begin{fichaepico}
\begin{campos}
  \item[Objetivo] O estudante acha a missão do dia, joga sem atrito e não perde
    o que fez --- em qualquer missão, antes mesmo de ela começar.
  \item[Para quem] Estudante, do 1º ano EF à 3ª série EM.
  \item[Valor] O que acontece \emph{dentro} de cada missão está no épico
    \id{E06}. Este épico é o que é igual em todas: a porta de entrada.
  \item[Estado] \pronto.
\end{campos}
\end{fichaepico}

\begin{tabularx}{\linewidth}{@{}L{16mm}X L{20mm}@{}}
\toprule
\headrow \thd{ID} & \thd{Feature} & \thd{Estado}\\
\midrule
\id{F05.1} & Minhas missões & \pronto\\
\zebra \id{F05.2} & Jogar a missão dentro da plataforma & \pronto\\
\bottomrule
\end{tabularx}

% ---------------------------------------------------------------------------
\feature{F05.1}{Minhas missões}{\pronto}

\begin{historia}{H05.1.1}{Ver as missões da minha turma}{\pronto}
\hu{Estudante}{ver as missões do meu ano e saber qual está aberta}{começar a
missão que o professor pediu}
\aceite
\begin{itemize}
  \item \enquote{Minhas missões} (\enquote{Olá, \{nome\}}) mostra a turma
    (ex.: \enquote{Turma 8A \textperiodcentered{} 8º ano \textperiodcentered{}
    manhã}) e um cartão por missão da grade do ano.
  \item Cartão aberto: \enquote{Começar}. Cartão fechado: cadeado e
    \enquote{Ainda não abriu}. Missão encerrada: etiqueta \enquote{encerrada}
    e \enquote{Esta você já fechou}.
  \item Aparecem só as missões da série da turma, mais as \enquote{para
    qualquer ano}.
  \item A missão fechada \textbf{não some}: o estudante vê que ela existe.
  \item Missão desligada pela Unooria para a rede some da grade.
\end{itemize}
\end{historia}

\begin{historia}{H05.1.2}{Entender quando não há nada para fazer}{\pronto}
\hu{Estudante}{uma mensagem clara quando não tenho missão}{não achar que a
plataforma quebrou}
\aceite
\begin{itemize}
  \item Mensagens diferentes para cada caso: \enquote{Você ainda não está em
    uma turma}, \enquote{Nenhuma missão por aqui ainda} e \enquote{Não
    conseguimos carregar suas missões}.
\end{itemize}
\end{historia}

% ---------------------------------------------------------------------------
\feature{F05.2}{Jogar a missão dentro da plataforma}{\pronto}

\begin{historia}{H05.2.1}{Abrir a missão}{\pronto}
\hu{Estudante}{abrir a missão com um toque}{começar a aula junto com a turma}
\aceite
\begin{itemize}
  \item A missão abre dentro de uma moldura da plataforma --- em formato de
    celular ou de tela cheia, conforme a missão --- com \enquote{\ensuremath{\leftarrow}{}
    Missões} no cabeçalho.
  \item Missão não liberada, de outra série ou desligada não abre, nem pelo
    endereço digitado.
  \item Para a criança, a mensagem de missão desligada é a mesma de missão
    ainda não aberta, para não quebrar a ficção.
\end{itemize}
\end{historia}

\begin{historia}{H05.2.2}{Não perder o que fiz}{\pronto}
\hu{Estudante}{que o que eu fiz na missão fique guardado}{não recomeçar se a
página recarregar ou a internet oscilar}
\aceite
\begin{itemize}
  \item O que a criança fez fica guardado na plataforma e volta depois de
    recarregar a página.
  \item Quando a janela da missão perde o foco, a missão pausa (os cronômetros
    congelam) e retoma ao voltar.
  \item A sessão não expira enquanto a criança está na missão (\id{H01.1.3}).
\end{itemize}
\pendencia{se a conexão cair de vez, os registros de uso ainda não enviados se
perdem (o progresso da missão, não). Guardar esses registros sem conexão é
evolução planejada.}
\end{historia}

\begin{historia}{H05.2.3}{Esperar o professor sem confusão}{\pronto}
\hu{Estudante}{saber quando estou esperando o professor}{não achar que travou}
\aceite
\begin{itemize}
  \item Várias missões param a turma até o professor soltar: a largada do
    relógio (unooGPT), a Aula 2 (Eco no Grupo, Rastro do Boato), o debate
    (Manifesto da IA), a roda de conversa (O Que a Foto Conta) e a revelação,
    a síntese e a reflexão (Quem mandou isso?).
  \item Nesses momentos a tela mostra uma espera clara, nunca um erro.
\end{itemize}
\end{historia}

\begin{historia}{H05.2.4}{Nunca ser comparado com os colegas}{\pronto}
\hu{família do estudante}{que nenhuma tela mostre nota ou ranking da
criança}{que a missão seja espaço de aprender, não de competir}
\aceite
\begin{itemize}
  \item Nenhuma tela do estudante mostra nota ou posição em ranking de
    colegas reais.
  \item Quando a missão tem ranking como parte da ficção (UnooTube), os
    colegas do ranking são fictícios.
  \item A correção acontece na plataforma; a resposta certa não chega ao
    aparelho antes da hora.
\end{itemize}
\end{historia}
